% Options for packages loaded elsewhere
\PassOptionsToPackage{unicode}{hyperref}
\PassOptionsToPackage{hyphens}{url}
\documentclass[
  11pt,
  a4paper]{article}
\usepackage{xcolor}
\usepackage{amsmath,amssymb}
\setcounter{secnumdepth}{5}
\usepackage{iftex}
\ifPDFTeX
  \usepackage[T1]{fontenc}
  \usepackage[utf8]{inputenc}
  \usepackage{textcomp} % provide euro and other symbols
\else % if luatex or xetex
  \usepackage{unicode-math} % this also loads fontspec
  \defaultfontfeatures{Scale=MatchLowercase}
  \defaultfontfeatures[\rmfamily]{Ligatures=TeX,Scale=1}
\fi
\usepackage{lmodern}
\ifPDFTeX\else
  % xetex/luatex font selection
  \setmainfont[]{Latin Modern Roman}
\fi
% Use upquote if available, for straight quotes in verbatim environments
\IfFileExists{upquote.sty}{\usepackage{upquote}}{}
\IfFileExists{microtype.sty}{% use microtype if available
  \usepackage[]{microtype}
  \UseMicrotypeSet[protrusion]{basicmath} % disable protrusion for tt fonts
}{}
\makeatletter
\@ifundefined{KOMAClassName}{% if non-KOMA class
  \IfFileExists{parskip.sty}{%
    \usepackage{parskip}
  }{% else
    \setlength{\parindent}{0pt}
    \setlength{\parskip}{6pt plus 2pt minus 1pt}}
}{% if KOMA class
  \KOMAoptions{parskip=half}}
\makeatother
\usepackage{longtable,booktabs,array}
\usepackage{calc} % for calculating minipage widths
% Correct order of tables after \paragraph or \subparagraph
\usepackage{etoolbox}
\makeatletter
\patchcmd\longtable{\par}{\if@noskipsec\mbox{}\fi\par}{}{}
\makeatother
% Allow footnotes in longtable head/foot
\IfFileExists{footnotehyper.sty}{\usepackage{footnotehyper}}{\usepackage{footnote}}
\makesavenoteenv{longtable}
\setlength{\emergencystretch}{3em} % prevent overfull lines
\providecommand{\tightlist}{%
  \setlength{\itemsep}{0pt}\setlength{\parskip}{0pt}}
\usepackage{amsmath,amssymb,amsthm}
\usepackage{booktabs}
\usepackage{graphicx}
\usepackage[margin=1in]{geometry}
\usepackage{hyperref}
\usepackage{xcolor}
\usepackage{unicode-math}
\hypersetup{colorlinks=true,linkcolor=blue!60!black,citecolor=green!50!black,urlcolor=blue!60!black}
\usepackage{fancyhdr}
\pagestyle{fancy}
\fancyhf{}
\fancyhead[L]{\small PACSeries Summary}
\fancyhead[R]{\small Dawn Field Institute}
\fancyfoot[C]{\thepage}
\renewcommand{\headrulewidth}{0.4pt}
\usepackage{bookmark}
\IfFileExists{xurl.sty}{\usepackage{xurl}}{} % add URL line breaks if available
\urlstyle{same}
\hypersetup{
  hidelinks,
  pdfcreator={LaTeX via pandoc}}

\title{PACSeries v0.2: Dawn Field Theory - A Summary}
\author{Peter Groom}
\date{February 2026}

\begin{document}
\maketitle

\renewcommand*\contentsname{Contents}
{
\setcounter{tocdepth}{3}
\tableofcontents
}
\section{PACSeries v0.2: Dawn Field Theory --- A
Summary}\label{pacseries-v0.2-dawn-field-theory-a-summary}

\textbf{Peter Groom}\\
Dawn Field Institute\\
February 2026

\begin{center}\rule{0.5\linewidth}{0.5pt}\end{center}

\subsection{Abstract}\label{abstract}

This document summarises the PACSeries v0.2, a collection of six
preprints that develop Dawn Field Theory from a single axiom --- PAC
(Potential-Actualization Conservation) --- through thermodynamic
foundations, mathematical structures, physical predictions, and
computational validation. The series establishes that information
erasure necessarily creates correlational structure, that the balance
constant governing recursive-conservation systems decomposes as
\(\Xi = \gamma + \ln\varphi \approx 1.0584\), and that this structural
framework produces closed-form expressions for the Feigenbaum constants
(to 13 significant figures), Standard Model parameters (to 5.7 ppm), and
the necessity of three spatial dimensions --- all from a conservation
law that also describes the internal dynamics of trained neural
networks. Each paper states its falsification conditions. This summary
provides the derivation chain, headline results, and reading guide for
the full series.

\begin{center}\rule{0.5\linewidth}{0.5pt}\end{center}

\subsection{1. The Central Idea}\label{the-central-idea}

Traditional physics treats information as something that
\emph{describes} pre-existing structure. Dawn Field Theory inverts this:
\textbf{information gradients may drive the emergence of structure
itself}.

The framework rests on one axiom:

\[f(\text{Parent}) = \sum f(\text{Children})\]

This is PAC conservation --- when potential becomes actual, the total is
conserved but redistributed. The unique stable solution to the two-term
PAC recursion \(\Psi(k) = \Psi(k{+}1) + \Psi(k{+}2)\) is
\(\Psi(k) = \varphi^{-k}\), where \(\varphi = (1{+}\sqrt{5})/2\) is the
golden ratio. This yields a natural information unit of
\(\ln\varphi \approx 0.4812\) per recursion level.

The six papers trace the consequences of this axiom across
thermodynamics, number theory, nonlinear dynamics, particle physics,
classical electromagnetism, and machine learning.

\begin{center}\rule{0.5\linewidth}{0.5pt}\end{center}

\subsection{2. The Derivation Chain}\label{the-derivation-chain}

\begin{verbatim}
AXIOM: PAC conservation — f(Parent) = Σ f(Children)
  → RECURSION: Ψ(k) = Ψ(k+1) + Ψ(k+2)
    → SOLUTION: Ψ(k) = φ^(-k)  [unique stable]
      → INFO UNIT: ΔI = ln(φ)
        → Paper 1: A/(A+ξ) → ln(φ)  [0.15% error, thermodynamic]
        → Paper 2: Ξ = γ + ln(φ)    [5 domains, p < 0.0003]
        → Paper 3: Feigenbaum from F₁₀ = 55  [13 digits]
        → Paper 4: sin²θ_W = 3/13, α to 5.7 ppm
        → Paper 5: Maxwell from depth-2 PAC, D = 3
        → Paper 6: PAC conservation in trained neural networks
\end{verbatim}

Each step derives from the previous. No step requires accepting any step
beyond it.

\begin{center}\rule{0.5\linewidth}{0.5pt}\end{center}

\subsection{3. Paper Summaries}\label{paper-summaries}

\subsubsection{Paper 1: The Structure Cost of
Erasure}\label{paper-1-the-structure-cost-of-erasure}

\textbf{Core claim}: Landauer's principle (erasure costs energy)
combined with the data processing inequality (DPI) \emph{requires} that
erasing information into a multi-mode environment creates new
correlational structure \(\xi\) between environmental modes. This
structure is topological --- invariant under temperature changes from
100K to 5000K --- not thermodynamic.

\textbf{Headline results}: - The collapse efficiency ratio
\(A/(A{+}\xi)\) converges toward \(\ln\varphi\) at 0.15\% proximity
(\(N = 5 \times 10^6\), Miller--Madow corrected) - Cascade coupling
self-sustains with 53× amplification over single events
(\(p = 2.75 \times 10^{-35}\)) - Temporal asymmetry: 69× early-vs-late
computational density (\(p = 3.25 \times 10^{-5}\))

\textbf{Significance}: This paper starts from two of the most secure
results in information theory (Landauer's principle and the DPI) and
shows that structure creation is \emph{mandatory}, not optional. The
golden ratio appears as a consequence of the conservation law, not as an
imposed parameter.

\subsubsection{Paper 2: The Balance Constant and Its
Decomposition}\label{paper-2-the-balance-constant-and-its-decomposition}

\textbf{Core claim}: The balance constant \(\Xi\) governing the boundary
between ordered and disordered computation decomposes as
\(\Xi = \gamma + \ln\varphi\), where \(\gamma \approx 0.5772\) is the
Euler--Mascheroni constant and \(\ln\varphi \approx 0.4812\). These
constants arise from \emph{different mathematics} --- \(\gamma\) from
the Mertens product over primes, \(\ln\varphi\) from PAC recursion ---
and converge independently.

\textbf{Headline results}: - Five independent domains (Fibonacci
arithmetic, cellular automata, prime sieve, Landauer erasure, Möbius
field dynamics) converge on \(\Xi \approx 1.058\) with \(p < 0.0003\) -
PAC conservation holds exactly across all 126 steps of the Eratosthenes
sieve (\(N = 500{,}000\)) - Class IV cellular automata cluster at
\(\Xi\) with \(p < 10^{-7}\) - Möbius field dynamics provides the
highest-precision measurement at 0.036\% from \(\gamma + \ln\varphi\),
from a continuous dynamical system on non-trivial topology

\textbf{Significance}: The convergence of five independent computational
substrates --- including both discrete (integers, automata) and
continuous (spectral fields on a Möbius manifold) --- is the strongest
evidence that \(\Xi\) is a universal constant rather than a
domain-specific artefact.

\subsubsection{Paper 3: Feigenbaum Constants from Fibonacci
Arithmetic}\label{paper-3-feigenbaum-constants-from-fibonacci-arithmetic}

\textbf{Core claim}: The Feigenbaum constants \(r_\infty\), \(\delta\),
and \(|\alpha|\) --- which have resisted closed-form expression for
nearly 50 years --- can be written as formulas using only \(\pi\),
Fibonacci numbers, and small integers. The integers are structurally
identified (\(55 = F_{10}\), \(52 = F_{10} - F_4\), etc.).

\textbf{Headline results}: - \(r_\infty\) matched to \textbf{13
significant figures} (relative error \(1.16 \times 10^{-14}\)) -
\(\delta\) matched to \textbf{8 significant figures}
(\(1.20 \times 10^{-9}\)) - \(|\alpha|\) matched to \textbf{6
significant figures} (\(4.02 \times 10^{-7}\)) - Probability of the best
triple occurring by chance: 1 in 280 billion - Exhaustive search of
3,920,499 parameter combinations finds only one triple achieving 7+
digit precision

\textbf{Significance}: This is arguably the hardest result in the series
to dismiss. No causal claim is made --- the paper reports the formulas,
proves their precision, and leaves the explanation open. If the
Feigenbaum constants genuinely have Fibonacci structure, the
implications for universality in nonlinear dynamics are significant.

\subsubsection{Paper 4: Standard Model Parameters from Fibonacci
Arithmetic}\label{paper-4-standard-model-parameters-from-fibonacci-arithmetic}

\textbf{Core claim}: A significant subset of the Standard Model's
parameters can be expressed as closed-form Fibonacci ratios from PAC
recursion. The gauge group \(U(1) \times SU(2) \times SU(3)\) is the
\emph{unique} combination whose adjoint dimensions (1, 3, 8) are all
Fibonacci numbers, closing at \(F_7 = 13\).

\textbf{Headline results}: - Gauge couplings to \textbf{5.7 ppm}; mixing
angles to 0.19\%; mass ratios to 5 ppm - Weak mixing angle
\(\sin^2\theta_W = 3/13\) matches running value at \(Q \approx M_W\) -
\(M_W/M_Z\) predicted at 0.03\% error - Koide formula for lepton mass
ratios to 0.5 ppm - Combined probability \(p < 10^{-5}\)

\textbf{Falsifiable predictions}: Z' boson at \(395 \pm 20\) GeV with
coupling \(g_{Z'}/g_Z = 1/13\); She--Lévêque turbulence constant
\(k(4) = 20\) in 4D simulations.

\textbf{Significance}: This is the most ambitious paper and will attract
the most scrutiny. The numerical matches are precise, but the
theoretical grounding for \emph{why} PAC recursion should constrain
gauge structure is interpretive, not derived. The falsifiable
predictions provide concrete experimental tests.

\subsubsection{Paper 5: Classical Physics from Information
Geometry}\label{paper-5-classical-physics-from-information-geometry}

\textbf{Core claim}: Maxwell's equations --- curl operations,
inverse-square forces, and three spatial dimensions --- emerge as
consequences of information-theoretic constraints (PAC, SEC, MED) rather
than independent empirical postulations.

\textbf{Headline results}: - Five independent arguments each require
exactly \(D = 3\) spatial dimensions: MED node bounds, curl algebra
closure, Möbius embedding, orbital stability, quaternion uniqueness -
SEC wave equation produces electromagnetic propagation at speed \(c\) -
Curl structure emerges from depth-2 recursion projection - Charge
quantisation: fractional quark charges (\(\pm 1/3\), \(\pm 2/3\)) follow
from MED nodes \(\leq 3\) - Gravity extension: Fibonacci depth
\(183 = F_7^2 + F_7 + 1\) gives \(F_{183} \approx 10^{38}\) (the
electromagnetic-to-gravitational hierarchy ratio)

\textbf{Honest negative}: The Mersenne--Fibonacci correspondence holds
for \(d = 1, 3, 7\) but fails at \(d = 15\).

\subsubsection{Paper 6: Computational Validation of PAC
Conservation}\label{paper-6-computational-validation-of-pac-conservation}

\textbf{Core claim}: PAC conservation --- discovered in thermodynamics
and validated across physics --- also describes the internal dynamics of
transformer-based language models without being explicitly programmed.
Hallucination is a direct PAC violation.

\textbf{Headline results}: - SEC phase classification predicts token
accuracy with zero free parameters: crystallised predictions achieve
100\% accuracy; chaotic 17--22\% (\(p < 0.0001\)) - \(\Xi\) appears in
trained weight spectra at 2.36× above random baselines
(\(\chi^2 = 5511\), \(p \approx 0\)) - Hallucination produces +9.6\%
uncompensated entropy - Enforcing PAC conservation (TinyCIMM-Boltzmann)
yields 16× reduction in context-switching shock with no measurable cost
to factual learning (\(p = 0.42\), n.s.)

\textbf{Honest negative}: \(\varphi\)-enrichment in top-2 ratios was
falsified --- identified as a softmax artefact.

\textbf{Significance}: The extension to neural networks provides a
practical application pathway and demonstrates that PAC conservation is
not merely a mathematical curiosity but a constraint that manifests in
trained computational systems.

\begin{center}\rule{0.5\linewidth}{0.5pt}\end{center}

\subsection{4. Cross-Cutting Themes}\label{cross-cutting-themes}

\subsubsection{What the series establishes
(measurement)}\label{what-the-series-establishes-measurement}

\begin{itemize}
\tightlist
\item
  Structure creation is mandatory for multi-mode erasure (Paper 1, from
  DPI + Landauer)
\item
  Five computational domains converge on \(\Xi \approx 1.058\)
  (\(p < 0.0003\), Paper 2)
\item
  Feigenbaum constants have Fibonacci expressions to 6--13 digits (Paper
  3)
\item
  PAC conservation holds exactly in the prime sieve (Paper 2)
\item
  SEC phase predicts transformer accuracy with zero free parameters
  (Paper 6)
\end{itemize}

\subsubsection{What the series derives
(analytical)}\label{what-the-series-derives-analytical}

\begin{itemize}
\tightlist
\item
  PAC forces emergence depth \(\leq 2\) (PAC→MED theorem, Paper 2 §9.3)
\item
  Only \(k = 2\) (Fibonacci) cascade produces \(\ln\varphi\) decay
  (Paper 2)
\item
  Curl structure requires and produces \(D = 3\) (Paper 5)
\item
  \(SU(2)\) and \(SU(3)\) are the only non-abelian gauge groups with
  Fibonacci adjoint dimensions (Paper 4)
\end{itemize}

\subsubsection{What the series proposes
(interpretation)}\label{what-the-series-proposes-interpretation}

\begin{itemize}
\tightlist
\item
  \(\gamma\) as ``discrete-to-continuous regularisation cost'' (Paper 2
  --- acknowledged: \(1/\sqrt{3}\) performs comparably)
\item
  Standard Model parameters from Fibonacci recursion (Paper 4 ---
  precise matches, interpretive grounding)
\item
  Gravity from SEC symmetric projection at Fibonacci depth 183 (Paper 5
  --- speculative)
\end{itemize}

\begin{center}\rule{0.5\linewidth}{0.5pt}\end{center}

\subsection{5. Falsification}\label{falsification}

Each paper states explicit falsification conditions. The series as a
whole would be undermined by:

\begin{enumerate}
\def\labelenumi{\arabic{enumi}.}
\tightlist
\item
  A new recursive-conservation system producing a balance constant
  \(> 1\%\) from \(\gamma + \ln\varphi\)
\item
  The Feigenbaum formulas being shown to arise from parameter fitting
  artefacts
\item
  Trained neural networks in architectures beyond Pythia/GPT-2 showing
  no \(\Xi\) enrichment
\item
  A construction of consistent vector electrodynamics in \(D \neq 3\)
\item
  The Z' prediction at \(395 \pm 20\) GeV being excluded by collider
  data
\end{enumerate}

The series prioritises precision, honest negatives, and clear separation
of measurement from interpretation. Failed predictions are documented
(phase-boundary narrative in Paper 2, \(\varphi\)-enrichment softmax
artefact in Paper 6, Mersenne--Fibonacci failure at \(d = 15\) in Paper
5).

\begin{center}\rule{0.5\linewidth}{0.5pt}\end{center}

\subsection{6. How to Read This Series}\label{how-to-read-this-series}

\textbf{If you have 10 minutes}: Read this summary, then Paper 3
(Feigenbaum). The 13-digit match is the single hardest result to
dismiss.

\textbf{If you have an hour}: Read Papers 1--3 in order. The derivation
chain from Landauer erasure through the balance constant to Feigenbaum
is self-contained and relies only on established mathematics.

\textbf{If you want to verify}: Each paper includes a complete
publication package --- Code (numbered experiment scripts), Data (JSON
results), and Figures (publication PNGs). Run
\texttt{python\ reproduce.py} in any paper's Code/ directory.

\textbf{If you want to falsify}: Check the falsification conditions in
each paper. Run the experiments with different parameters. The code is
open.

\begin{center}\rule{0.5\linewidth}{0.5pt}\end{center}

\subsection{7. Publication Details}\label{publication-details}

\begin{longtable}[]{@{}
  >{\raggedright\arraybackslash}p{(\linewidth - 8\tabcolsep) * \real{0.1458}}
  >{\raggedright\arraybackslash}p{(\linewidth - 8\tabcolsep) * \real{0.1458}}
  >{\raggedright\arraybackslash}p{(\linewidth - 8\tabcolsep) * \real{0.1875}}
  >{\raggedright\arraybackslash}p{(\linewidth - 8\tabcolsep) * \real{0.2708}}
  >{\raggedright\arraybackslash}p{(\linewidth - 8\tabcolsep) * \real{0.2500}}@{}}
\toprule\noalign{}
\begin{minipage}[b]{\linewidth}\raggedright
Paper
\end{minipage} & \begin{minipage}[b]{\linewidth}\raggedright
Title
\end{minipage} & \begin{minipage}[b]{\linewidth}\raggedright
Version
\end{minipage} & \begin{minipage}[b]{\linewidth}\raggedright
Experiments
\end{minipage} & \begin{minipage}[b]{\linewidth}\raggedright
Key Metric
\end{minipage} \\
\midrule\noalign{}
\endhead
\bottomrule\noalign{}
\endlastfoot
1 & The Structure Cost of Erasure & 2.1 & 19 &
\(A/(A{+}\xi) \to \ln\varphi\) at 0.15\% \\
2 & The Balance Constant & 2.2 & 15 & 5 domains, \(p < 0.0003\) \\
3 & Feigenbaum Constants & 2.1 & 9 & \(r_\infty\) to 13 digits \\
4 & Standard Model Parameters & 2.1 & 14 & \(\alpha\) to 5.7 ppm \\
5 & Classical Physics & 2.1 & 9 & 5 independent \(D{=}3\) arguments \\
6 & Computational Validation & 2.1 & 10 & \(\Xi\) at 2.36× in weight
spectra \\
\end{longtable}

\textbf{Total}: 76 numbered experiments, 72 data files, 37 figures.

\textbf{Repository}:
\href{https://github.com/dawnfield-institute/dawn-field-theory}{github.com/dawnfield-institute/dawn-field-theory}\\
\textbf{Prior version}: PACSeries v0.1 (October 2025), Zenodo DOI:
\href{https://zenodo.org/records/17295103}{10.5281/zenodo.17295103}

\begin{center}\rule{0.5\linewidth}{0.5pt}\end{center}

\subsection{References}\label{references}

Each paper contains its own reference list. Cross-references use the
notation ``Paper N, §M'' throughout the series. All code dependencies
are standard scientific Python (NumPy, SciPy, Matplotlib).

\begin{center}\rule{0.5\linewidth}{0.5pt}\end{center}

\emph{PACSeries v0.2. All code, data, and figures are publicly available
under AGPL-3.0 (code) and CC-BY-4.0 (papers).}

\end{document}
